\documentclass[11pt,a4paper]{article}

\usepackage[utf8]{inputenc}
\usepackage[spanish]{babel}
\usepackage[T1]{fontenc}
\usepackage{lmodern}
\usepackage[left=2.5cm,right=2.5cm,top=3cm,bottom=3cm]{geometry}
\usepackage{graphicx}
\usepackage{xcolor}
\usepackage{fancyhdr}
\usepackage{titlesec}
\usepackage{hyperref}

% Configuración de colores
\definecolor{primarycolor}{RGB}{0, 71, 143} % Azul médico / corporativo
\definecolor{secondarycolor}{RGB}{43, 43, 43}

% Configuración de hipervínculos
\hypersetup{
    colorlinks=true,
    linkcolor=primarycolor,
    urlcolor=primarycolor,
    pdftitle={Manual de Usuario: Staff Interno},
    pdfauthor={RadiographXpress}
}

% Formato de encabezado y pie de página
\pagestyle{fancy}
\fancyhf{}
\fancyhead[L]{\textcolor{secondarycolor}{\small \textbf{RadiographXpress}}}
\fancyhead[R]{\textcolor{secondarycolor}{\small Manual de Staff Interno}}
\fancyfoot[C]{\thepage}
\renewcommand{\headrulewidth}{0.4pt}
\renewcommand{\footrulewidth}{0.4pt}

% Formato de títulos
\titleformat{\section}{\Large\bfseries\color{primarycolor}}{\thesection}{1em}{}
\titleformat{\subsection}{\large\bfseries\color{secondarycolor}}{\thesubsection}{1em}{}

% Directorio de imágenes
\graphicspath{{screenshots/}}

\title{
    \vspace{-2cm}
    \includegraphics[width=0.4\textwidth]{logo_placeholder.png} \\ % Reemplazar con el logo real si existe en Overleaf
    \vspace{1cm}
    \Huge \textbf{RadiographXpress} \\
    \Large \vspace{0.5cm} Manual de Usuario y Operaciones: \\ Staff Interno y Panel Principal
}
\author{Personal Administrativo, Asistentes y Radiólogos}
\date{\today}

\begin{document}

\maketitle
\thispagestyle{empty}

\begin{abstract}
\noindent Este manual está dirigido al personal que labora y opera directamente con el sistema interno de RadiographXpress. Contiene los lineamientos técnicos, flujos de trabajo e instrucciones visuales paso a paso para Administradores de Sistema, Asistentes de mostrador y Médicos Reportadores (Radiólogos).
\end{abstract}

\newpage
\tableofcontents
\newpage

%%%%%%%%%%%%%%%%%%%%%%%%%%%%%%%%%%%%%%%%%%%%%%%%%%%%%%%%%%%%%%%%%%%%
% PARTE 1 - ADMINISTRADOR DEL SISTEMA
%%%%%%%%%%%%%%%%%%%%%%%%%%%%%%%%%%%%%%%%%%%%%%%%%%%%%%%%%%%%%%%%%%%%
\section{Administrador del Sistema (Django Admin)}

Este perfil está dirigido al Administrador de Negocios o Administrador Técnico. El panel de administración proporciona control total sobre los permisos, usuarios, roles y todos los datos en la base de datos subyacente de RadiographXpress.

\subsection{Acceso al Panel de Control Administrativo}
El panel de control es un portal oculto y de acceso restringido diseñado exclusivamente para el personal gerencial y de TI.

\begin{enumerate}
    \item Navega a la URL de administración del sistema (generalmente \texttt{http://[tudominio.com]/admin}).
    \item Ingresa tus credenciales de superusuario o de staff autorizadas.
    \item Al ingresar, serás dirigido al Panel Principal donde se listan todas las entidades ("Aplicaciones") de la base de datos de RadiographXpress.
\end{enumerate}

\begin{figure}[h]
    \centering
    \includegraphics[width=0.9\textwidth]{admin_dashboard.png}
    \caption{Vista Principal del Panel de Administración.}
\end{figure}

\subsection{Gestión de Registros y Roles de Usuario}
El núcleo del sistema RadiographXpress se basa en \textbf{Grupos} de permisos (\texttt{Groups}). 

\subsubsection{Usuarios (\texttt{Users})}
\begin{enumerate}
    \item En la sección \textbf{Authentication and Authorization}, haz clic en \textbf{Users}.
    \item Aquí verás todos los usuarios registrados en el sistema, sea cual sea su rol.
    \item \textbf{Modificar un usuario:} Si un usuario olvida su contraseña o necesita ser dado de baja, haz clic en su correo electrónico (o `username`).
    \begin{itemize}
        \item \textbf{Cambiar contraseña:} Existe un formulario especializado para cambiar claves de acceso directamente.
        \item \textbf{Activar/Desactivar:} Desmarcando la casilla \texttt{Active} puedes suspender temporalmente el acceso de un usuario.
    \end{itemize}
\end{enumerate}

\begin{figure}[h]
    \centering
    \includegraphics[width=0.9\textwidth]{admin_users.png}
    \caption{Gestión general de usuarios del sistema.}
\end{figure}

\subsubsection{Administración de Perfiles Específicos}
\textbf{Médicos Reportadores (\texttt{doctorsDashboard}):}
Si contratas a un radiólogo nuevo, deberás crearle un usuario y asignarlo a este perfil, añadiendo además su Firma Digital.

\vspace{0.5cm}
\noindent \textbf{Cuidado de los Datos:} Siempre desactiva usuarios (\texttt{Active=False}) en lugar de eliminarlos para mantener la integridad forense y médico-legal del historial.

\newpage

%%%%%%%%%%%%%%%%%%%%%%%%%%%%%%%%%%%%%%%%%%%%%%%%%%%%%%%%%%%%%%%%%%%%
% PARTE 2 - ASISTENTE
%%%%%%%%%%%%%%%%%%%%%%%%%%%%%%%%%%%%%%%%%%%%%%%%%%%%%%%%%%%%%%%%%%%%
\section{Asistente / Administrativo (Mostrador)}

Como asistente del sistema, eres primordial en el orden de los accesos y de la carga inicial de los pacientes antes de su entrada al RIS/PACS.

\subsection{Inicio de Sesión y Dashboard}
Inicia sesión con las credenciales proporcionadas. Tendrás herramientas de validación de identidad y alta de perfiles.

\begin{figure}[h]
    \centering
    \includegraphics[width=0.9\textwidth]{assistant_dashboard.png}
    \caption{Panel principal (Dashboard) del Asistente.}
\end{figure}

\subsection{Gestión Manual de Pacientes}
\begin{enumerate}
    \item Utiliza la barra superior para buscar pacientes globales por nombre o registro, evitando duplicados.
    \item \textbf{Creación de Nuevos Pacientes:} Si llega un paciente a realizarse un estudio y no tiene cuenta, haz clic en el formulario pertinente y captura demográficos completos y precisos.
    \item Al agendar un estudio, este se vincula creando un \texttt{Study Request}.
\end{enumerate}

\subsection{Validación y Aprobación de Médicos Asociados}
Para médicos que soliciten darse de alta como Referentes (Asociados), **tú tienes la llave final**. Debes verificar sus licencias antes de aprobarlos en el sistema para que puedan ver los resultados de sus pacientes.
\begin{itemize}
    \item Ve al listado de \textit{Pendientes de Aprobación}.
    \item Verifica la Cédula Profesional en los portales gubernamentales u oficiales locales.
    \item Haz clic en \textbf{Aceptar} (Vincular). El médico será notificado vía email inmediatamente.
\end{itemize}

\vspace{0.3cm}
\noindent \textit{(Nota: Las capturas relacionadas con los correos enviados se encuentran detalladas en el archivo de anexos SMTP y configuraciones adicionales de sistema).}

\newpage

%%%%%%%%%%%%%%%%%%%%%%%%%%%%%%%%%%%%%%%%%%%%%%%%%%%%%%%%%%%%%%%%%%%%
% PARTE 3 - MÉDICO REPORTADOR (RADIÓLOGO)
%%%%%%%%%%%%%%%%%%%%%%%%%%%%%%%%%%%%%%%%%%%%%%%%%%%%%%%%%%%%%%%%%%%%
\section{Médico de Interpretación (Radiólogo)}

Este es el portal de operaciones de diagnóstico. Como Médico Reportador, aquí dictaminarás, clasificarás y enviarás los resultados a pacientes y médicos tratantes a través de nuestro PACS.

\subsection{Ingreso y Tu Tablero de Trabajo (Worklist)}
Al ingresar a tu cuenta de radiólogo, visualizarás todas las solicitudes de la clínica en tiempo real. 

\begin{figure}[h]
    \centering
    \includegraphics[width=0.9\textwidth]{doctor_dashboard.png}
    \caption{Dashboard Médico - Worklist y estado de estudios.}
\end{figure}

\subsection{Firma Digital (Obligatoria)}
Al configurar tu perfil por primera vez desde las opciones de \textbf{Mi Perfil}, carga el escaneo o archivo digital de tu firma (\texttt{.png} transparente recomendado). Se inyectará automáticamente en los documentos.

\subsection{Bloqueo e Interpretación de un Estudio}
Debido a que pueden existir múltiples radiólogos trabajando a la vez, se utilizan bloqueos (\textit{Locks}).

\begin{enumerate}
    \item Ubica un estudio con estado \textbf{"Pending"} y ábrelo.
    \item A partir de ese segundo, ningún otro radiólogo podrá tomar el estudio (\textit{Bloqueo Inteligente}).
    \item Redacta la nota médica estructurada: \textit{Descripción clínica, Hallazgos, Conclusiones y Recomendaciones}.
    \item Una vez presiones Enviar/Finalizar, el sistema firmará y despachará el PDF instantáneamente. El estado pasará a \textbf{"Completado"}.
\end{enumerate}

\end{document}
